% !TeX root = main.tex
\lecture{16}{Thu 12 Mar 2026 13:00}{Statistical Physics III: Atmosphere and Boltzmann Factors}

Recap from last lecture:
\begin{itemize}
    \item We derived:
    \[
        \rho_N(h) = \rho_N(0) \exp \left(- \frac{mgh}{k_B T}\right)
    \]
    \[
        P(h) = P(0) \exp \left(- \frac{mgh}{k_B T}\right)
    \]
    
    To give the number density and pressure as an exponential function of height in the atmosphere.

    \item Since $mgh / k_B T$ is being passed into an exponential, it must be dimensionless. Hence $k_B T / mg$ must have dimensions of $h$. 
    \item This leads to the scale height $h_0$ where $P(h_0) = P(0) / e $
\end{itemize}


\section{Atmosphere Continued}

We have established that the number density (number of fluid molecules per unit volume), at height 0, $\rho_N(0)$ can be related to atmospheric pressure ($P_\text{at} = P(0)$) by:
\[
    \rho_N(0) = \frac{P_\text{atmospheric}}{k_B T}
\]

To determine the total number of gas molecules in the slab of atmosphere with cross-sectional area $A$, we can integrate number density across all $h$:
\begin{align*}
    N &= A \int_{0}^{\infty} \rho_N(h) \, dh \\
    &= A \int_{0}^{\infty} \rho_N(0) \exp\left(- \frac{mgh}{k_B T}\right) \, dh \\
    &= A \rho_N(0) \int_{0}^{\infty} \exp\left(- \frac{mgh}{k_B T}\right) \, dh \\
    &= A \rho_N(0) \left[- \frac{k_B T}{mg} \exp \left(- \frac{mhg}{k_B T}\right)\right]_0^\infty \\
    &= -A \rho_N(0) \frac{k_B T}{mg} \left[0 - 1\right] \\
    N &= A \rho_N(0) \frac{k_B T}{mg} \\
    \implies \rho_N(0) &= \frac{Nmg}{A k_B T}
\end{align*}

And using:
\[
    h_0 = \frac{k_B T}{mg}
\]
\[
    \rho_N(0) = \frac{N}{A h_0}
\]

This has a physical definition, being the average number density of the slab of atmosphere between $h = 0, h = h_0$

We further define a new quantity:
\[
    \Pr(h) = \frac{A \rho_N(h)}{N}
\]

Which gives the probability of (for a particle) finding it at height $h$

\subsection{Probability Normalisation}
We start from:
\[
    Pr(h) = \frac{A \rho_N(h)}{N}
\]
\[
    h_0 = \frac{k_B T}{mg}
\]
\[
    \rho_N(h) = \rho_N(0) \exp(-mgh / k_B T)
\]

Hence:
\[
    \Pr(h) = \frac{A}{N} \times \frac{N}{A h_0} \exp(-mhg / k_B T)
\]
\[
    \Pr(h) = \frac{mg}{k_B T} \exp(-mhg / k_B T)
\]



\section{Boltzmann Factors}
\textcolor{red}{This is on the exam.} We have the quantity $mgk / k_B T$, which is dimensionless. What does this quantity actually mean? $mgh$ is a term for the potential energy (gravitational in this example), $k_B T$ is the thermal energy
